% v11 Threats (TV1-TV7 plus retained v10 threats) and Discussion (DI1-DI5).
\section{Threats to Validity}\label{sec:threats}

\subsection{Construct Validity}
Every outcome is a property of locations and labels, not of review concerns. A correct start line does not show that a comment still applies, and a changed label does not show that a concern is affected. The controlled edit mix is uniform over twelve operators chosen by us, and one of them, re-indent, often yields invalid or behaviour-changing Python (Section~\ref{sec:stress}). Automatically sampled anchors need not resemble the spans supervisors select: supervisors may favour short, repeated or boundary-spanning lines, which the stress suites show to be harder. Changes outside the quoted range can matter to a concern and are not evaluated. The effort model's weights are illustrative, not measured costs.

\subsection{Oracle Validity}
The real-history reference is a syntax-tree correspondence procedure, not semantic ground truth: it aligns statements by normalized tree structure and uses heuristics for replacements, moves and deletions, and does not decide whether two statements behave alike. It excludes \rExcludedPct\% of candidate observations; excluded files have larger diffs (median \wpaExclDiffMedian{} changed lines versus \wpaRetainedDiffMedian{} for retained ones), so accuracy is conditional on a decidable and easier subset and may overestimate performance on unrestricted code evolution. The original manual audit covered only 40 decisions and was carried out by members of the development team, which risks confirmation bias; a 70-case blinded packet with separately held answers is prepared, but no independent responses have been collected. The human-annotated external benchmark of Section~\ref{sec:external} is the only reference in this paper made independently of the authors. Every target the reference calls unchanged keeps its quoted lines at the reference location, but for relocated targets that location can be a copy that already existed, so its identification errors can overstate missed re-checks and also hide a few: they bound the deployed miss rate between \refdirBoundPct\% and \refcopyMissPctWorst\%. Placement accuracy and unnecessary re-checks have no such bound.

\subsection{Internal Validity}
The default threshold was selected on a separate synthetic seed. Leave-one-project-out calibration, the factorial ablation and the guards' real-history results use histories that earlier revisions had inspected, so they are retrospective; the guards' held-out evidence is limited to stress suites generated after their protocol was fixed, the external benchmark and the student-authored repositories. The corrected copy guard AGi was specified after the first guard's student and compound results had been read, so its held-out evidence is the fresh stress seed of Section~\ref{sec:locator}; CodeMapper with the cross-file fallback and AGi was composed after outcomes were read and is exploratory. The guarded ten-commit replay was first run with the rejected provenance rule and repeated with AGi under a recorded amendment before its outcomes existed. The pinned CodeMapper is not fully deterministic: recomputed on the same chains, it changed its answer on \wiChIGateChains{} of 4,000, so CodeMapper rates carry a small run-to-run variation beyond their sampling intervals. The protocols are timestamped locally, not externally preregistered. The conventional baselines are reimplementations of documented behaviours, not tests of the deployed products, and implementation details of GitHub, Gerrit or GitLab may differ.

\subsection{Statistical Dependence}
Observations from the same project, commit pair or source file are correlated. Wilson intervals are nominal per-observation intervals. We therefore also report cluster-bootstrap intervals by source file, commit pair and project; one internal comparison loses its resolved direction under project clustering (Section~\ref{sec:design-results}). With nine projects, project-level intervals are coarse. The paper reports many paired comparisons and does not adjust them for multiplicity. The registered hypotheses are therefore stated per population and contrast together with the number that held (H1 was supported on \wiHypOneSupported{} of \wiHypOneTotal{} contrasts and H2 on \wiHypTwoSupported{} of \wiHypTwoTotal{}), and an individual interval that excludes zero should be read as one comparison among many, not as a confirmed effect.

\subsection{External Validity}
All real histories are Python, from four internal projects and five widely used libraries; none consists of supervised student resubmissions. Student revisions may restructure more, repeat more code, be submitted incompletely or omit files. The public student-authored repositories of Section~\ref{sec:student} are a stand-in that shares authorship but not the supervision setting. The algorithm operates on text, but cross-language effectiveness rests on the external benchmark of Section~\ref{sec:external}, not on our own corpora. Errors may accumulate over many resubmissions; our multi-hop analysis replays real commit chains without independent intermediate identity labels or supervisor intervention.

\subsection{Implemented but Not Evaluated}
Media-note reliability for users, privacy and access control and concurrent composition are implemented design controls. Section~\ref{sec:statecheck} reports the scenario tests that exercise them, which are pass/fail checks of specified behaviour, not a security assessment, penetration test or user study.

\section{Discussion}\label{sec:discussion}

\subsection{What Was and Was Not Shown}
The experiments show where anchored code is found, how its change is labelled and how often the label withholds a re-check, on defined distributions. They do not show that supervisors begin a re-review faster, revisit more of their earlier points, close fewer issues in error or that students respond better. Those benefits motivate the design and are plausible, but they are hypotheses. The direct test is a within-subjects study in which supervisors re-review the same revised submissions with a conventional interface and with \sys{}, measuring time to locate earlier comments, the share of earlier issues correctly revisited, erroneous closures and workload.

\subsection{Matcher Strength}
The method is strongest where its design assumptions hold: code that is shifted, moved or edited lightly and that has no identical copies nearby. On naturally edited statements its default threshold is too strict, the earlier v1 design relocates more, an independent contextual line tracker relocates more again within files (Sections~\ref{sec:realtracker} and~\ref{sec:student}), and a dedicated region mapper, CodeMapper, locates human-annotated regions far more accurately and recovers their boundaries where the cascade mostly does not (Section~\ref{sec:external}). The cascade is not a competitive locator. Its distinctive capabilities are cross-file recovery and the labels that drive the issue state. Because the labels are computed from the evidence a matcher returns, the relocation engine is replaceable. Section~\ref{sec:locator} did this with CodeMapper: the same label rules and window guards (IG and AGi; the variable-height window is part of the cascade's own search) made it safer on every population without changing its placement. The workflow's safety properties are those of the combined system, and a deployment should pair the strongest available locator with the safety layer rather than choose between them.

\subsection{Costs of Different Failures}
The failures do not cost the same. A lost anchor is visible: the issue requests re-check and the supervisor searches for the code, a cost of time. An unnecessary re-check costs a short confirmation. A missed re-check or a silent misplacement costs attention that is never prompted, and can lead a supervisor to judge a concern against the wrong code. The evaluation therefore treats the last two as safety failures and the first two as effort, and recommends guards that convert silent failures into visible ones even at the price of more confirmations.

Whether that price is worth paying depends on how many carried issues point at code that changed, which no study here measured. Let $p$ be that share and $w$ the cost of a missed re-check in units of one unnecessary confirmation. The recommended guards (Table~\ref{tab:locator}) cost no more than the deployed cascade when $w\,p$ times the fall in the missed re-check rate is at least $(1-p)$ times the rise in the false re-check rate. With $w=1$ the break-even share is \wiRealBreakEvenOne\% on real histories, \wiStuBreakEvenOne\% on the student-authored repositories and \wiHumBreakEvenOne\% on the human-annotated regions; with $w=10$ it is \wiRealBreakEvenTen\%, \wiStuBreakEvenTen\% and \wiHumBreakEvenTen\%. These points ignore silent misplacements, which the guards also reduce on every population, so they favor the deployed cascade. For comparison, \wpeOutdatedPct\% of review threads in our open-source sample were outdated by GitHub, an upper bound on changed code rather than a measured share of unresolved concerns. A course that weighs a hidden change at ten confirmations or more gains from the guards on all three populations once more than \wiHumBreakEvenTen\% of its carried issues point at changed code.

\subsection{Misplacement as a Traceability Problem}
A misplaced anchor attaches feedback to code with a different meaning. The interface shows the label in plain language beside every relocated anchor, and the supervisor can edit an anchor; every edit is an attributed revision. When tier~3 accepts a \lbl{changed} window, however, the carried anchor's quotation is replaced by the matched text, and later versions are matched against that updated quotation. In the deployed code the original quotation survives only in the previous version's issue, reachable through the lineage reference. The multi-hop replay shows the consequence: after a wrong fuzzy update, later revisions can match the updated quotation exactly, so the latest label alone cannot certify continuity with the original target. The identity rule of Section~\ref{sec:identity} keeps the reviewed quotation in every carried anchor and, in replay, made the latest request carry the whole history (Section~\ref{sec:prefixes}). It is evaluated in replay only; the deployed portal does not yet store it.

\subsection{Media Notes}
Voice, video and image notes follow their issue or are re-found with the cascade when pinned to code. Their placement reliability is that of code anchors, and their handling across versions is covered by the scenario tests of Section~\ref{sec:statecheck}. No quantitative relocation study or user study of media notes was performed.

\subsection{Deploying the Method in a Course}
An institution adopting the method should keep the local search scope and tier order, choose the threshold as a policy between relocation and honest removal (Section~\ref{sec:design-results}), enable the indentation guard and the corrected copy guard AGi over whichever locator it uses (Section~\ref{sec:locator}), keep the reviewed quotation beside the current one, prefer few anchors per issue, warn supervisors when they anchor a recurring one-line span, and keep resolution with the supervisor. Supervisors' corrections of misplaced anchors should be logged, with consent, so that a course can calibrate locally.
